\section{Conclusion}

Agents are becoming a new class of internet user. Google AX shows how to run them as isolated, scalable workloads, but a sandbox only protects the cluster from the agent; it does nothing to help the rest of the world decide whether to trust what the agent asks for. Know Your Agent offers the vocabulary for that decision (identity, principal, authorization and behaviour) but, today, no open reference implementation tied to a real agent platform.

AgentPassport joins the two. It keeps AX unchanged and adds a thin, standards-based layer around it: every task gets a short-lived passport naming the human behind it and the exact permissions granted; every request leaving the sandbox is signed by a proxy the agent can neither influence nor bypass; every service can verify those requests with a few lines of code; and a trust observer stops agents whose behaviour stops making sense. The worked examples show the intended outcome: a prompt-injected agent that tries to move 9,000~TRY is blocked twice and suspended within seconds, a looping agent is stopped before it burns the budget, and a stolen passport is worthless.

The project's contribution is therefore both practical and conceptual: working open-source components that fill gaps AX itself lists on its roadmap, and a concrete, measurable answer to the question \emph{``how should an agent introduce itself, and how should the world check?''}
